% ============================================================
% EJERCICIO 6 — Conversión de imagen a LaTeX (Mathpix / LaTeX-OCR)
% Tareas: 1) Elija una ecuación (foto de su cuaderno o de un libro).
%         2) Conviértala con Mathpix o LaTeX-OCR.
%         3) Pegue el código ORIGINAL de la herramienta y su versión corregida.
% Entregue: imagen, este .tex y PDF.
% ============================================================
\documentclass[12pt,a4paper]{article}
\usepackage[utf8]{inputenc}
\usepackage[T1]{fontenc}
\usepackage[spanish]{babel}
\usepackage{amsmath,amssymb}
\usepackage{booktabs}
\usepackage{siunitx}
\sisetup{locale=FR, output-decimal-marker={,}}
\usepackage{circuitikz}
\usepackage[hidelinks]{hyperref}
\usepackage{graphicx}
\usepackage{fancyvrb}

\begin{document}
\section{Ecuación original}
% TODO \includegraphics[width=0.6\textwidth]{ecuacion.png}  (ponga su imagen)
\fbox{\parbox{0.8\textwidth}{\centering\vspace{1cm}Imagen de la ecuación\vspace{1cm}}}

\section{Código generado por la herramienta}
% Herramienta usada: TODO
\begin{Verbatim}[frame=single,fontsize=\small]
TODO pegue aquí el LaTeX tal como lo devolvió la herramienta
\end{Verbatim}

\section{Código corregido y resultado}
% TODO ¿Qué errores tenía? (símbolos, \left/\right, alineación…)
\begin{equation}
  % TODO su versión corregida
\end{equation}

\section{Reflexión}
% TODO 3-4 líneas: ¿qué tan precisa fue la herramienta? ¿qué falló?
\end{document}
% RÚBRICA (20 pts): ecuación compila (5) · ecuación fiel al original (5)
% · código original y corregido mostrados (5) · reflexión crítica (5)
